\section{Linear Divergence of Correlations in QCD Vacuum  }\label{sec:otherresult}
%{\it Results:}\ 

In this section, we analyze the linear divergence of 
the Wilson line in the matrix elements taken in the QCD
vacuum. We try to test
if the divergence is universal and if the divergent mass can be extracted successfully from these matrix elements. We consider
three new types of matrix elements: 1) large-size Wilson loop
that has been used to extract heavy-quark potential, 2) the
vacuum matrix element of the quasi-PDF operator, 3) the vacuum matrix element of Wilson link operator in a fixed gauge. 
The data are for MILC ensembles and, when applicable, with the clover action.
According to the last section, smearing does not change
the form of renormalization, and therefore, we consider
here the higher-precision data with one step smearing. Since we just want to test the linear divergence, we can use the simple function Eq.~(\ref{eq:logM_ss}) in our fitting. 

\subsection{Wilson Loop}

To extract the linear divergence from lattice data, 
a Wilson loop is a good choice because it is easy to calculate, 
gauge invariant, and has high precision. It has been extensively studied in
the literature~\cite{Chen:2016fxx,Zhang:2017bzy,Musch:2010ka,Green:2017xeu,Chen:2017gck}, particularly in relation to the heavy quark potential. 
%In Fig.~\ref{fig:vev2}, we have shown the heavy quark potential as a function of lattice spacing. 

\begin{figure}[tbp]
\centering
\includegraphics[width=8cm]{images/fitting_Wilson_Loop_hyp1.pdf}
\includegraphics[width=8cm]{images/ert_vs_r+t_Wilson_Loop_hyp1.pdf}
\caption{Test of linear divergence for a Wilson Loop for HYP smearing cases.
The slope is consistent approximately with hadronic matrix elements from the previous section.}
\label{fig:LP_test}
\end{figure}

In our analysis, we will not consider 
the heavy quark potential interpretation but simply consider
a Wilson loop as a matrix element and fit
its $1/a$ divergence by looking at the logarithm
of the matrix element, as shown in 
Fig.~\ref{fig:LP_test}. The lattice
spacing dependence has been fitted well
with our formula with a choice of $\Lambda_{\rm QCD}=
0.39$ GeV, which is a value favored by
our fitting in the previous section. After isolating the linear
divergent coefficient, we plot it as a function
of $r+t$, for which half of the slope gives us
the divergent coefficient. Our value is
about 2.5 GeV$^{-1}$fm$^{-1}$ (0.49 if dimensionless). This value is very similar
to the value we found in the previous section (Fig.~\ref{fig:e(z)_hyp}),
demonstrating that the linear divergence 
can be very well extracted from the Wilson loop.
 
\subsection{Vacuum Matrix Element of Quasi-PDF Operator}

\begin{figure}[tbp]
\centering
\includegraphics[width=8cm]{images/fitting_VEV_hyp1_Oa.pdf}
\includegraphics[width=8cm]{images/ez_vs_z_VEV_hyp1_Oa.pdf}
\caption{Test of linear divergence for the vacuum matrix of the quasi-PDF operator. The matrix
element is much suppressed and there is no leading twist contribution.}
\label{fig:vev12}
\end{figure}

Here we consider the vacuum matrix elements of a quasi-PDF operator
which were suggested as choices for the renormalization
factor in Ref.~\cite{Braun:2018brg}. 

One can consider the vacuum expectation value (VEV) of $O_{\gamma_t}(t)$. As a gauge invariant choice proposed in Ref.~\cite{Braun:2018brg}. $\langle O_{\Gamma}(t)\rangle$
can be obtained through the following stochastic estimation:
\bea\label{eq:vev1}
&&\langle O_{\Gamma}(t)\rangle =\frac{1}{L_s^3}\langle \sum_{\vec{x}}\textrm{Tr}[U(\vec{x},0;\vec{x},t)\Gamma S_w(\vec{x},t)]\rangle\nonumber\\
&=&\frac{1}{L_s^3}\langle \sum_{\vec{x}}\textrm{Tr}[U(\vec{x},0;\vec{x},t)\Gamma S(\vec{x},t;\vec{x},0)]\rangle\nonumber\\
&&+\frac{1}{L_s^3}\sum_{\vec{x},\vec{y},\vec{y}\neq\vec{x}}\langle \textrm{Tr}[U(\vec{x},0;\vec{x},t)\Gamma S(\vec{x},t;\vec{y},0)]\rangle
\eea
where the second term on the right hand side of the second line vanishes as it is not gauge invariant, and $S_w(\vec{x},t)=\sum_{\vec{y}}S(\vec{x},t;\vec{y},0)$ is a wall source quark propagator without gauge fixing. Such a proposal can only be applied for the non-vanishing cases with $\Gamma=\gamma_t$ or ${\cal I}$. We will apply it to $O_{\gamma_t}$.  

However, it is simple to see that the $O_{\gamma_t}$ matrix element in the vacuum vanishes at small $t$ using operator product expansion.
Therefore the matrix element is susceptible 
to large O($a$) correction. Only at very large $t$,
one might see the correct slope. 

Results of a test of the linear divergence for the vacuum matrix element of the PDF operator is shown in Fig.~\ref{fig:vev12}. 
The slope at large $t$ appears to be consistent with that from the Wilson loop case but only within uninterestingly large errors.

\subsection{Landau-gauge-fixed Wilson link}

For the gauge-dependent matrix elements, we consider 
the Wilson line in Landau gauge as the simplest choice. 

\begin{figure}[tbp]
\centering
\includegraphics[width=8cm]{images/fitting_Wilson_Link_hyp1.pdf}
\includegraphics[width=8cm]{images/ez_vs_z_Wilson_Link_hyp1.pdf}
\caption{Test of linear divergence for Landau gauge fixed Wilson link for HYP smearing cases. At small-$z$, it can be calculated in perturbation theory where the linear divergence works well. For large-$z$, the matrix element does not have a transfer-matrix 
interpretation in this gauge.}
\label{fig:link_test}
\end{figure}

The result of our test of the linear divergence is shown in Fig.~\ref{fig:link_test}.
At small-$z$, where perturbation theory works well, the slope roughly agrees with 
the prediction from  perturbation theory. For large-$z$, the matrix element does not
have a transfer-matrix interpretation in this gauge.

\begin{figure}[tbp]
\centering
\includegraphics[width=8cm]{images/ez_vs_z_Comp_Oa.pdf}
\caption{Comparison of linear divergences in different vacuum matrix elements.}
\label{fig:com_test}
\end{figure}

Finally, Fig.~\ref{fig:com_test} is a comparison of the linearly divergent term for Wilson Loop, VEV and Wilson link. At small $z$, the linear divergence terms of Wilson Loop and Wilson link are consistent while at large $z$ they disagree strongly. The linear divergence term of VEV is different from the others because the large discretization errors 
make the extraction of the linear divergence very imprecise.
